\documentclass[12pt]{article}
%^ Start of document
\usepackage{times}  %Makes text TimesNewRoman
\usepackage{setspace} %Allows you to set single or double space 
\usepackage{biblatex} %Imports biblatex package
\usepackage{graphicx} % Required for inserting images
\usepackage{authblk}
\usepackage[colorlinks=true, urlcolor=blue, linkcolor=blue]{hyperref}
\usepackage{subfig}
\usepackage{float}
\usepackage{tabularx}
\usepackage{multirow}
\usepackage[paperheight=11in,
   paperwidth=8.5in,
   top=20mm,
   bottom=20mm,
   left=40mm,
   right=40mm]{geometry}
\usepackage{moresize}
\usepackage{tikz}
\usepackage{xcolor}
\usepackage{physics}
\usepackage{amssymb}
\usepackage{mathrsfs}
\usepackage{amsmath}
\usepackage{ragged2e}
\usepackage{comment}
\usepackage{xcolor}
\addbibresource{ElectroHW3.bib}
\begin{document}
\singlespacing  %Sets single spacing for whole document
\title{Classical Electrodynamics HW 3}

\author[1]{Zachary Tullier}
\affil[1]{Student\\Physics Department\\ University of Houston - Clear Lake \\Houston, Texas\\U.S}

\maketitle


\section{Textbook Question 3.6} \label{section1}
    \underline{Question : }
    Two point charges $q$ and $-q$ are located on the z axis at $z = +a$ and $z = -a$, respectively.
    \begin{enumerate}
        \item [a)] Find the electrostatic potential as an expansion in spherical harmonics and powers of $r$ for both $r>a$ and $r<a$.
        \item [b)] Keeping the product $qa \equiv \frac{p}{2}$ constant, take the limit of $a \rightarrow{} 0$ and find the potential for $r \neq 0$. This is by definition a dipole along the $z$ axis and it's potential. 
    \end{enumerate}

    %\newpage
    \underline{Question 3.6 Answer a): }
    Integrate equation \ref{eq: a2-12} from Appendix \ref{Appendix2}. The charges are located on the z axis so two of the axis have no component.
    \begin{equation} \label{eq: 3.6-15}
        \Phi = \frac{1}{4 \pi \epsilon_0} \left(\frac{q}{\abs{\vec{r} - a \hat{k}}} - \abs{\vec{r} + a \hat{k}}\right)
    \end{equation}
    Use the addition theorem where $r_<$ is smaller of ($r,a$), and $r_>$ is larger of ($r,a$)
    \begin{equation} \label{eq: 3.6-16}
        \frac{1}{\abs{\vec{r}-\vec{r}'}} = 4 \pi \sum_{l=0}^\infty \sum_{m=-l}^l \frac{1}{2l+1} \frac{r_<^l}{r_>^{l+1}} Y_{lm}(\theta', \phi') Y_{lm}(\theta, \phi) 
    \end{equation}
    Apply addition theorem to equation, using $\phi'=0$ for both cases, 
    \newline and $\theta'=\begin{cases} 0 \text{for $r-a\hat{k}$} \\ \pi \text{for $r+a\hat{k}$} \end{cases}$ because the charges are on the same axis.
    \begin{equation} \label{eq: 3.6-17}
        \Phi = \frac{q}{4 \pi \epsilon_0} \left(4 \pi \sum_{l=0}^\infty \sum_{m=-l}^l \frac{1}{2l+1} \frac{r_<^l}{r_>^{l+1}} Y_{lm}^*(0, 0) Y_{lm}(\theta, \phi) -4 \pi \sum_{l=0}^\infty \sum_{m=-l}^l \frac{1}{2l+1} \frac{r_<^l}{r_>^{l+1}} Y_{lm}^*(\pi, 0) Y_{lm}(\theta, \phi)  \right)
    \end{equation}
    Simplify
    \begin{equation} \label{eq: 3.6-18}
        \Phi = \frac{q}{\epsilon_0} \sum_{l=0}^\infty \sum_{m=-l}^l \frac{1}{2l+1} \frac{r_<^l}{r_>^{l+1}} Y_{lm}(\theta, \phi) \left( Y_{lm}^*(0, 0) - Y_{lm}^*(\pi, 0) \right)
    \end{equation}
    \begin{equation} \label{eq: 3.6-19}
        \Phi = 
        \begin{cases}
            \frac{q}{\epsilon_0} \sum_{l=0}^\infty \sum_{m=-l}^l \frac{1}{2l+1} \frac{r_l}{a^{l+1}} Y_{lm}(\theta, \phi) \left( Y_{lm}^*(0, 0) - Y_{lm}^*(\pi, 0) \right) \text{when $r<a$}
            \\
            \frac{q}{\epsilon_0} \sum_{l=0}^\infty \sum_{m=-l}^l \frac{1}{2l+1} \frac{a^l}{r^{l+1}} Y_{lm}(\theta, \phi) \left( Y_{lm}^*(0, 0) - Y_{lm}^*(\pi, 0) \right) \text{when $r>a$}
        \end{cases}
    \end{equation}

    \underline{Question 3.6 Answer b): }
    Simplify equation by applying azimuthally symmetric condition, all terms other than $m=0$ vanish.
    \begin{equation} \label{eq: 3.6-20}
        \Phi = \frac{q}{\epsilon_0} \sum_{l=0}^\infty  \frac{1}{2l+1} \frac{r_<^l}{r_>^{l+1}} Y_{lm}(\theta, \phi) \left( Y_{lm}^*(0, 0) - Y_{lm}^*(\pi, 0) \right)
    \end{equation}
    Convert Spherical harmonics into Legendre polynomial
    \begin{equation} \label{eq: 3.6-21}
        \Phi = \frac{q}{4 \pi \epsilon_0} \sum_{l=0}^\infty   \frac{r_<^l}{r_>^{l+1}} P_{l}(cos(\theta)) \left(P_{l}(1) - P_{l}(-1) \right)
    \end{equation}
    Evaluate Legendre polynomial
    \begin{equation} \label{eq: 3.6-22}
        \Phi = \frac{q}{4 \pi \epsilon_0} \sum_{l=0}^\infty   \frac{r_<^l}{r_>^{l+1}} P_{l}(cos(\theta)) \left(1-{(-1)}^{-1} \right)
    \end{equation}
    Simplify and alter iteration
    \begin{equation} \label{eq: 3.6-23}
        \Phi = \frac{2q}{4 \pi \epsilon_0} \sum_{l=0, odd}^\infty   \frac{r_<^l}{r_>^{l+1}} P_{l}(cos(\theta)) 
    \end{equation}
    Take the limit of $\alpha \rightarrow{} 0$ and apply $r \neq 0$ and $r>a$
    \begin{equation} \label{eq: 3.6-24}
        \Phi = \frac{2q}{4 \pi \epsilon_0} \sum_{l=0, odd}^\infty   \frac{a^l}{r^{l+1}} P_{l}(cos(\theta)) 
    \end{equation}
    \begin{equation} \label{eq: 3.6-25}
        \Phi = \frac{p}{4 \pi \epsilon_0} \sum_{l=0, odd}^\infty   \frac{a^{l-1}}{r^{l+1}} P_{l}(cos(\theta)) 
    \end{equation}
    \begin{equation} \label{eq: 3.6-26}
        \Phi = \frac{p}{4 \pi \epsilon_0} \left(\frac{1}{r^2} P_1(cos(\theta)) + \frac{a^2}{r^4} P_3(cos(\theta)) + \frac{a^4}{r^6} P_5(cos(\theta)) +... \right) 
    \end{equation}
    Take limit $\alpha \rightarrow{} 0$
    \begin{equation} \label{eq: 3.6-27}
        \Phi = \frac{p}{4 \pi \epsilon_0} \frac{1}{r^2} cos(\theta)
    \end{equation}
    
    \underline{Question 3.6 Answer c): }
    Then if dipole is surrounded by a grounded sphere shell sharing an origin, there will be additional potential from the sphere
    \newline Using equation \ref{eq: 3.6-27}, add potential from sphere. The potential from the surface of a sphere can be derived from the Laplace equation in spherical coordinates (From appendix 2, section \ref{Appendix2}, specifically equation \ref{eq: a2-43}
    \begin{gather*}
        \Phi(r,\theta, \phi) = \sum_{l=0}^\infty (A_l r^l + B_l r^{-l-1}) P_l(cos(\theta))
    \end{gather*}
    Where $B_l=0$ because the region of valid solution includes the origin.
    \begin{equation} \label{eq: eq: 2-13.1}
        \Phi(r,\theta, \phi) = \sum_{l=0}^\infty A_l r^l P_l(cos(\theta))
    \end{equation}
    Add potential from sphere into equation \ref{eq: 3.6-27}
    \begin{equation} \label{eq: 2-14}
        \Phi = \frac{p}{4 \pi \epsilon_0} \frac{1}{r^2} cos(\theta) + \sum_{l=0}^\infty A_l r^l P_l(cos(\theta))
    \end{equation}
    Apply boundary condition; potential is zero at surface of sphere. Therefore, $\Phi=0$ if $r=b$ where $b$ is the radius of the sphere. 
    \begin{equation} \label{eq: 3.6-29}
        0 = \frac{p}{4 \pi \epsilon_0} \frac{1}{b^2} cos(\theta) + \sum_{l=0}^\infty A_l b^l P_l(cos(\theta))
    \end{equation}
    Only $l=1$ will be nonzero because of orthogonality
    \begin{equation} \label{eq: 3.6-30}
        0 = \frac{p}{4 \pi \epsilon_0} \frac{1}{b^2} cos(\theta) + A_l b cos(\theta)
    \end{equation}
    Isolate $A_l$
    \begin{equation} \label{eq: 3.6-31}
        A_l = -\frac{p}{4 \pi \epsilon_0} \frac{1}{b^3}
    \end{equation}
    Plug $A_l$ into equation \ref{eq: 2-14}
    \begin{equation} \label{eq: 3.6-32}
        \Phi = \frac{p}{4 \pi \epsilon_0} \frac{1}{r^2} cos(\theta) + \sum_{l=0}^\infty -\frac{p}{4 \pi \epsilon_0} \frac{1}{b^3} r^l P_l(cos(\theta))
    \end{equation}        
    Again, Only $l=1$ will be nonzero because of orthogonality
    \begin{equation} \label{eq: 3.6-33}
        \Phi = \frac{p cos(\theta)}{4 \pi \epsilon_0 b^2} \left(\left(\frac{b}{r} \right)^2 - \frac{r}{b} \right)
    \end{equation}
    
\newpage
\section{Textbook Question 3.7} \label{section2}
    \underline{Question: }
    Three point charges ($q, -2q, q$) are located in a straight line with separation $a$ and with the middle charge ($-2q$) at the origin of a grounded conducting spherical shell of radius $b$, as indicated in the sketch. 
    \begin{enumerate}
        \item [a)] Write down the potential of the three charges in the absence of the grounded sphere. Find the limit form of the potential as $a \rightarrow{} 0$, but the product $qa^2 = Q$ remains finite. Write this latter answer in spherical coordinates.
        \item [b)] The presence of the grounded sphere of radius $b$ alters the potential for $r<b$. The added potential can be viewed as caused by the surface-charge density induced on the inner surface at $r=b$ or by image charges located at $r>b$. Use linear superposition to satisfy the boundary conditions and find the potential everywhere inside the sphere for $r<a$ and $r>a$. Show that in the limit $a \rightarrow{} 0$, 
        \begin{gather*}
            \Phi(r,\theta, \phi) \rightarrow{} \frac{Q}{2 \pi \epsilon_0 r^3} \left(1 - \frac{r^5}{}b^5 \right) P_2(cos(\theta))
        \end{gather*}
    \end{enumerate}
    
    %\newpage
    \underline{Question 3.7, Section a) Answer: }
    Integrate equation \ref{eq: a2-12} from \ref{Appendix2} for all three charges
    \begin{equation} \label{eq: 3.7-1}
        \Phi = \frac{1}{4 \pi \epsilon_0} \left(\frac{q}{\abs{\vec{r} - \alpha(\hat{z})}} + \frac{q}{\abs{\vec{r}-\alpha(-\hat{z})}} - \frac{2q}{\abs{\vec{r}}} \right)
    \end{equation}
    Expand the first and second term using the Legendre Polynomial Expansion (Shown in Appendix 1, pg \pageref{Appendix1}) 
    \newline If $r>a$
    \begin{equation} \label{eq: 3.7-2}
        \Phi = \frac{1}{4 \pi \epsilon_0} \left(q \sum_{l=0}^\infty \frac{a_l}{r^{l+1}} P_l(cos(\theta)) + q - \sum_{l=0}^\infty \frac{{(-1)}^l a^l}{r^{l+1}} P_l(cos(\theta)) \frac{2q}{r} \right)        
    \end{equation}
    simplify
    \begin{equation} \label{eq: 3.7-3}
        \Phi = \frac{q}{4 \pi \epsilon_0} \left(\frac{-2}{r} \sum_{l=0}^\infty \left(1+{(-1)}^l \right)\frac{a_l}{r^{l+1}} P_l(cos(\theta)) \right)      
    \end{equation}
    Simplify by altering iteration
    \begin{equation} \label{eq: 3.7-4}
        \Phi = \frac{q}{2 \pi \epsilon_0} \left(\sum_{l=0, even}^\infty \frac{a_l}{r^{l+1}} P_l(cos(\theta)) \right)
    \end{equation}
    \newline If $r>a$
    \begin{equation} \label{eq: 3.7-5}
        \Phi = \frac{1}{4 \pi \epsilon_0} \left(q \sum_{l=0}^\infty \frac{r_l}{a^{l+1}} P_l(cos(\theta)) + q \sum_{l=0}^\infty \frac{{(-1)}^l r^l}{a^{l+1}} P_l(cos(\theta)) - \frac{2q}{r} \right)        
    \end{equation}
    simplify
    \begin{equation} \label{eq: 3.7-6}
        \Phi = \frac{q}{4 \pi \epsilon_0} \left(\frac{-2}{r} \sum_{l=0}^\infty \left(1+{(-1)}^l \right)\frac{a_l}{r^{l+1}} P_l(cos(\theta)) \right)      
    \end{equation}
    Simplify by altering iteration
    \begin{equation} \label{eq: 3.7-7}
        \Phi = \frac{q}{2 \pi \epsilon_0} \left(\sum_{l=0, even}^\infty \frac{r_l}{a^{l+1}} \left(1-\frac{a}{r} \delta_{l,0} \right) P_l(cos(\theta)) \right)      
    \end{equation}
    As $a \rightarrow{} 0$, the only potential that matters is the $r>a$, and we only keep the first non-zero term
    \newline set $qa^2 = Q$
    \begin{equation} \label{eq: 3.7-8}
        \Phi(r) = \frac{Q}{2 \pi \epsilon_0} \left(\frac{1}{r^3} P_2(cos(\theta)) \right)
    \end{equation}
    Evaluate Legendre polynomial
    \begin{equation} \label{eq: 3.7-9}
        \Phi(r) = \frac{Q}{4 \pi \epsilon_0 r^3} (3 cos^2(\theta) - 1)
    \end{equation}
    
    \underline{Question 3.7, Section b) Answer: }
    Integrate equation \ref{eq: a2-12} from Appendix 3, section \ref{Appendix2} for all three charges. In this integral we can place image charges azimuthally symmetric. We can begin with equation \ref{eq: 3.7-4} and add a element for the surface charge induced by the image charges.
    \begin{equation} \label{eq: 3.7-10}
        \Phi = \sum_l A_l r^l P_l(cos(\theta)) + \frac{q}{2 \pi \epsilon_0} \left(\sum_{l=0, even}^\infty \frac{a_l}{r^{l+1}} P_l(cos(\theta)) \right)
    \end{equation}
    Apply boundary condition $\Phi=0$ at surface ($r=b$)
    \begin{equation} \label{eq: 3.7-11}
        0 = \sum_l A_l r^l P_l(cos(\theta)) + \frac{q}{2 \pi \epsilon_0} \left(\sum_{l=0, even}^\infty \frac{a_l}{r^{l+1}} P_l(cos(\theta)) \right)
    \end{equation}
    Simplify
    \begin{equation} \label{eq: 3.7-12}
        0 =  A_l b^l + \frac{q}{2 \pi \epsilon_0} \frac{a^l}{b^{l+1}}
    \end{equation}
    Isolate $A_l$
    \begin{equation} \label{eq: 3.7-13}
        A_l =  -\frac{q}{2 \pi \epsilon_0} \frac{a^l}{b^{2l+1}}
    \end{equation}
    Substitute back into equation
    \begin{equation} \label{eq: 3.7-14}
        \Phi(r) = \sum_{l=2, even}^\infty \frac{q}{2 \pi \epsilon_0} \left(\frac{a}{b} \right)^l \left(\frac{b^l}{r^{l+1}} - \frac{r^l}{b^{l+1}} \right) P_l(cos(\theta))
    \end{equation}
    Since we are taking limit $a \rightarrow{} 0$, the $r<a$ case does not matter
    \newline Set $qa^2 = Q$, and we only need the first non-zero term
    \begin{equation} \label{eq: 3.7-15}
        \Phi(r) = \frac{Q}{2 \pi \epsilon_0 r^3} \left(1-\frac{r^5}{b^5} \right) P_2(cos(\theta))
    \end{equation}
    
    
\newpage
\section{Textbook Question 3.14} \label{section3}
    \underline{Question: }
    A line charge of length 2d with a total charge $Q$ has a linear charge density varying as $(d^2 - z^2)$, where $z$ is the distance from the midpoint. A grounded, conducting, spherical shell of inner radius $b>d$ is centered at the midpoint of the line charge. 
    \begin{enumerate}
        \item [a)] Find the potential everywhere inside the spherical shell as an expansion in Legendre polynomials 
        \item [b)] Calculate the surface-charge density induced on the shell.
        \item [c)] Discuss your answers to parts a and b in the limit that $d \ll b$.
    \end{enumerate}

    \underline{Question 3.14, section a) Answer: }
    Use the Green function for the Dirichlet boundary 
    \begin{equation} \label{eq: 3.14-1}
        \Phi(x) = \frac{1}{4 \pi \epsilon_0} \int \rho(x') G_D d^3x' - \frac{1}{4 \pi} \oint (\Phi \frac{\partial G_D}{\partial n'} da'
    \end{equation}
    Apply boundary condition $\Phi(x) = 0$ when $x=2d$
    \begin{equation} \label{eq: 3.14-2}
        \Phi(x) = \frac{1}{4 \pi \epsilon_0} \int \rho(x') G_D d^3x'
    \end{equation}
    If $r<d$ then
    \begin{equation} \label{eq: 3.14-3}
        \rho(r,\theta,\phi) = \frac{3 Q}{8 \pi d^3} \frac{(d^2 - r^2)}{r^2} \left(\delta(cos(\theta)-1) + \delta(cos(\theta)+1) \right)
    \end{equation}
    If $r>d$ then $\rho(r,\theta,\phi) = 0$
    \newline Substitute this back into equation
    \begin{equation} \label{eq: 3.14-4}
        \Phi(x) = \frac{1}{4 \pi \epsilon_0} \int_0^{2\pi} \int_0^\pi \int_0^\infty \rho(x') G_D {r'}^2 sin(\theta') dr' d\theta' d\phi'
    \end{equation}
    Integrate
    \begin{equation} \label{eq: 3.14-5}
        \Phi(x) = \frac{3 Q}{16 \pi \epsilon_0 d^3} \left(\int_0^d (d^2-{r'}^2) G_D(\theta'=0) dr' + \int_0^d (d^2 - {r'}^2 G_D(\theta'=\pi) dr' \right)
    \end{equation}
%    \newline \textcolor{red}{Lots of inbetween}
    \begin{equation} \label{eq: 3.14-6}
        G_D = 4 \pi \sum_{l=0}^\infty \sum_{m=-l}^l \frac{1}{2l+1} \left(\frac{r_<^l}{r_>^{l+1}} - \frac{(rr')^l}{b^{2l+1}} \right) Y_{lm}^*(\theta',\phi') Y_{lm}(\theta,\phi)
    \end{equation}
    Since there is azimuthal symmetry, all terms are removed except $m=0$
    \begin{equation} \label{eq: 3.14-7}
        G_D = 4 \pi \sum_{l=0}^\infty \left(\frac{r_<^l}{r_>^{l+1}} - \frac{(rr')^l}{b^{2l+1}} \right) P_l(cos(\theta') P_l(cos(\theta))
    \end{equation}
    Apply boundary conditions at $\theta'=0$
    \begin{equation} \label{eq: 3.14-8}
        G_D(\theta'=0) = \sum_{l=0}^\infty \left(\frac{r_<^l}{r_>^{l+1}} - \frac{(rr')^l}{b^{2l+1}} \right) P_l(cos(\theta))
    \end{equation}
    Apply boundary conditions at $\theta'=\pi$
    \begin{equation} \label{eq: 3.14-9}
        G_D(\theta'=\pi) = \sum_{l=0}^\infty \left(\frac{r_<^l}{r_>^{l+1}} - \frac{(rr')^l}{b^{2l+1}} \right) (-1)^l P_l(cos(\theta))
    \end{equation}
    Substitute back into equation
    \begin{equation} \label{eq: 3.14-10}
        \Phi(x) = \frac{3 Q}{8 \pi \epsilon_0 d^3} \sum_{l=0, even}^\infty P_l(cos(\theta)) \int_0^d dr'(d^2-{r'}^2) \left(\frac{r_<^l}{r_>^{l+1}} - \frac{(rr')^l}{b^{2l+1}} \right) 
    \end{equation}
    If $r>d$ then
    \begin{equation} \label{eq: 3.14-11}
        \Phi(x) = \frac{3 Q}{8 \pi \epsilon_0 d^3} \sum_{l=0, even}^\infty P_l(cos(\theta)) \int_0^d dr'(d^2-{r'}^2) \left(\frac{{r'}^l}{r^{l+1}} - \frac{(rr')^l}{b^{2l+1}} \right) 
    \end{equation}
    Evaluate integral
%    \newline \textcolor{red}{Do inbetween}
    \begin{equation} \label{eq: 3.14-12}
        \Phi(x) = \frac{3 Q}{8 \pi \epsilon_0 d^3} \sum_{l=0, even}^\infty P_l(cos(\theta)) \frac{d^l}{(l+3)(l+3)b^{l+1}} \left(\left(\frac{b}{r} \right)^{l+1} - \left( \frac{r}{b} \right)^l \right) 
    \end{equation}
    If $r<d$ then
    \begin{equation} \label{eq: 3.14-13}
        \begin{split}
        \Phi(x) = \frac{3 Q}{8 \pi \epsilon_0 d^3} \sum_{l=0, even}^\infty P_l(cos(\theta)) 
        \\\left( \int_0^r dr'(d^2-{r'}^2) \left(\frac{{r'}^l}{r^{l+1}} - \frac{(rr')^l}{b^{2l+1}} \right) + \int_0^d dr'(d^2-{r'}^2) \left(\frac{{r}^l}{{r'}^{l+1}} - \frac{(rr')^l}{b^{2l+1}} \right) \right)
        \end{split}        
    \end{equation}
    Evaluate Integral
%    \newline \textcolor{red}{Do inbetween}
    \begin{equation} \label{eq: 3.14-14}
        \begin{split}
        \Phi(x) = \frac{3 Q}{8 \pi \epsilon_0 d^3} \sum_{l=0, even}^\infty P_l(cos(\theta)) 
        \\ \left(\frac{d^2 (2l+1)}{l(l+1)} + \frac{r^2 (2l+1)}{(2-l)(3+l)} + r^l \left(\frac{2}{d^{l-2} l (l-2)} - \frac{2 d^{l+3}}{b^{2l+1} (l+3) (l+1)} \right) \right)
        \end{split}
    \end{equation}
    
    \underline{Question 3.14, section b) Answer: }
    Use the definition of the surface-charge density
%    \newline \textcolor{red}{Find how this value was formulated}
    \begin{equation} \label{eq: 3.14-15}
        \sigma = \left(-\epsilon_0 \frac{\partial \Phi}{\partial n} \right)_S
    \end{equation}
    Evaluate at the inner surface $r=b$ and inward $n=-r$
    \begin{equation} \label{eq: 3.14-16}
        \sigma = \left(\epsilon_0 \frac{\partial \Phi}{\partial r} \right)_{r=b}
    \end{equation}
    Substitute equation in
    \begin{equation} \label{eq: 3.14-17}
        \sigma = \left(\epsilon_0 \frac{\partial }{\partial r} \left( \frac{3 Q}{8 \pi \epsilon_0 d^3} \sum_{l=0, even}^\infty P_l(cos(\theta)) \frac{d^l}{(l+3)(l+3)b^{l+1}} \left(\left(\frac{b}{r} \right)^{l+1} - \left( \frac{r}{b} \right)^l \right) \right) \right)_{r=b}
    \end{equation}
    Derive and simplify
%    \newline \textcolor{red}{Do inbetween}
    \begin{equation} \label{eq: 3.14-18}
        \sigma = -\frac{3 Q}{4 \pi b^2} \sum_{l=0, even}^\infty P_l(cos(\theta)) \frac{(2l+1)}{(l+1)(l+3)} \left(\frac{d}{b} \right)^l
    \end{equation}
    
    \underline{Question 3.14, section c) Answer: }
    If $d \ll b$ then the case to evaluate is equation equation \ref{eq: 3.14-11}.
    \newline If $d \ll b$ then $\left(\frac{d}{b} \right)^l = 0$ except at $l=0$.
    \begin{equation} \label{eq: 3.14-19}
        \Phi(x) = \frac{Q}{4 \pi \epsilon_0} \left(\frac{1}{r} - \frac{1}{b} \right)
    \end{equation}
    Using equation \ref{eq: 3.14-19} and the definition of surface-charge density \ref{eq: 3.14-15}
%    \newline \textcolor{red}{Do inbetween}
    \begin{equation} \label{eq: 3.14-20}
        \sigma = -\frac{Q}{4 \pi b^2}
    \end{equation}
    

\section{Appendix} 
    \subsection{Legendre Polynomial Expansion} \label{Appendix1}
        \begin{equation} \label{eq: a1-1}
            \frac{1}{\abs{\vec{r}-\vec{r_0}}} = 
            \begin{cases}
                \sum_{l=0}^\infty \frac{r_0^l}{r^{l+1}} P_l(cos(\theta)) \textbf{  if $r>r_0$}
                \\
                \sum_{l=0}^\infty \frac{r^l}{r_0^{l+1}} P_l(cos(\theta)) \textbf{  if $r>r_0$}
            \end{cases}
        \end{equation}
        
    \subsection{Formulation of Electric Potential} \label{Appendix2}
        Two Electric point charges exert a force on each other.
        \begin{equation} \label{eq: a2-1}
            \vec{F} = q \vec{E} = k q_1 q_2 \frac{x_1-x_2}{\abs{x_1-x_2}^3}
        \end{equation}
        \newline Electric field at point x due to a point charge
        \begin{equation} \label{eq: a2-2}
            \vec{E} = k q_1 \frac{x-x_1}{\abs{x-x_1}^3}
        \end{equation}
        \newline In electrostatic units (esu) k = 1, In SI:
        \begin{equation} \label{eq: a2-3}
            k = \frac{1}{4 \pi \epsilon_0} = 10^{-7} c^2
        \end{equation}
        \newline Since Electrodynamics is a macroscopic theory, we will consider the "point" to be many charged sources. Experiment shows that the fields add linearly.
        \begin{equation} \label{eq: a2-4}
            \vec{E}(x)= k \sum_i q_i \frac{x-x_i}{\abs{x - x_i}^3}
        \end{equation}
        \newline Consider the many charged sources to be "small" (Replace $q$ with $\rho$) and "numerous" (Replace Sum with Integral)at point $x'$. Coulomb's law in terms of the field for a charge density. Replace singular x dimension with 3D x,y,z dimension 
        \begin{equation} \label{eq: a2-5}
            \vec{E}(x)= k \int \rho_q(x') \frac{x-x'}{\abs{x - x'}^3}d^3x'
        \end{equation}
        A vector field can be specified completely if its divergence and curl are given everywhere in space, except up to the gradient of scalar function that satisfies the Laplace equation. 
%        \newline \textcolor{red}{when you get to section 1.9, this Laplace stuff might come into play}
        \newline Therefore we need to specify the curl of $\vec{E}$
        \newline Starting at Coulomb's generalized integral form 
        \begin{gather*}
            \vec{E(x)}= k \int \rho_q(x^{'}) \frac{x-x^{'}}{\abs{x - x'}^3}d^3x'
        \end{gather*}
%        \newline \textcolor{red}{I guess I don't know how $\nabla$ works. Figure out in between}
        \newline Using Property of Nabla ($\nabla$)
        \begin{equation} \label{eq: a2-6}
            -\nabla(\frac{1}{\abs{x-x'}}) = \frac{x-x'}{\abs{x - x'}^3}
        \end{equation}
        Substitute equation \ref{eq: a2-3}
        \begin{equation} \label{eq: a2-7}
            \vec{E(x)}= k \int \rho_q(x') -\nabla(\frac{1}{\abs{x-x'}}d^3x'
        \end{equation}
        Nabla ($\nabla$) is constant; therefore,
        \begin{equation} \label{eq: a2-8}
            \vec{E(x)}= -k \nabla \int \rho_q(x') (\frac{1}{\abs{x-x'}}d^3x'
        \end{equation}
        Cross nabla ($\nabla$) on both sides of equation and substitute equation
        \begin{equation} \label{eq: a2-9}
            \nabla \cross \vec{E(x)}= \nabla \cross k -k\nabla \int \rho_q(x') (\frac{1}{\abs{x-x'}}d^3x'
        \end{equation}
        Therefore,
        \begin{equation} \label{eq: a2-10}
            \nabla \cross \vec{E} = 0
        \end{equation}
        By manipulating equation
        \begin{equation} \label{eq: a2-10}
            \vec{E(x)}= \nabla \int (\frac{k \rho_q (x')}{\abs{x-x'}}d^3x'
        \end{equation}
        We can now define a portion of the equation as the scalar potential ($\Phi$)
        \begin{equation} \label{eq: a2-11}
            \vec{E(x)}= \nabla \Phi
        \end{equation}
        Therefore the scalar potential is defined as 
        \begin{equation} \label{eq: a2-12}
            \Phi = \int (\frac{k \rho_q (x')}{\abs{x-x'}}d^3x'
        \end{equation}

    \subsection{Laplace Equation in Spherical Coordinates} \label{Appendix2}
        The definition of the Laplace equation
        \begin{equation} \label{eq: a2-1}
            \nabla^2 \Phi = 0
        \end{equation}
        The definition of $\nabla$ in spherical coordinates
        \begin{equation} \label{eq: a2-2}
            \nabla = \hat{e}_1 \frac{\partial}{\partial r} + \hat{e}_2 \frac{1}{\theta} \frac{\partial}{\partial \theta} + \hat{e}_3 \frac{1}{r sin(\theta)}\frac{\partial}{\partial \phi}
        \end{equation}
        Where $r$ is the radial distance from the origin, $\theta$ is the polar angle with the z-axis, $\phi$ is the azimuthal angle in the x-y plan relative to the x-axis
        \newline Combine the two
        \begin{equation} \label{eq: a2-3}
            \frac{1}{r} \frac{\partial^2}{\partial r^2} (r \Phi) + \frac{1}{r^2 sin(\theta)} \frac{\partial}{\partial \theta} \left(sin(\theta) \frac{\partial \Phi}{\partial \theta} \right) + \frac{1}{r^2 sin^2(\theta)} \frac{\partial^2 \Phi}{\partial \phi^2} = 0
        \end{equation}
        
        Use method of separation of variables by trying a solution of the form
        \begin{equation} \label{eq: a2-4}
            \Phi(r,\theta,\phi) = \frac{R(r)}{r} P(\theta) Q(\phi)
        \end{equation}
        Substitute equation \ref{eq: a2-4} into equation \ref{eq: a2-3}
        \begin{equation} \label{eq: a2-5}
            \frac{1}{r} \frac{\partial^2}{\partial r^2} (r \frac{R}{r} PQ + \frac{1}{r^2 sin(\theta)} \frac{\partial}{\partial \theta} (sin(\theta) \frac{\partial \frac{R}{r} PQ}{\partial \theta}) + \frac{1}{r^2 sin^2(\theta)} \frac{\partial^2 \frac{R}{r} PQ}{\partial \phi^2} = 0
        \end{equation}
        Multiply by $\frac{r^3 sin^2(\theta)}{RPQ}$
        \begin{equation} \label{eq: a2-6}
            \frac{r^2 sin^2(\theta)}{R} \frac{d^2 R}{dr^2} + \frac{sin(\theta)}{P} \frac{d}{d\theta}(sin(\theta) \frac{dP}{d\theta}) + \frac{1}{Q} \frac{d^2Q}{d\phi^2} = 0
        \end{equation}
        Isolate the last term
        \begin{equation} \label{eq: a2-7}
            \frac{r^2 sin^2(\theta)}{R} \frac{d^2 R}{dr^2} + \frac{sin(\theta)}{P} \frac{d}{d\theta}(sin(\theta) \frac{dP}{d\theta}) = - \frac{1}{Q} \frac{d^2Q}{d\phi^2}
        \end{equation}
        Equate it to a constant ($m$)
        \begin{equation} \label{eq: a2-7}
            \frac{1}{Q(\phi)} \frac{d^2Q(\phi)}{d\phi^2} = -m^2
        \end{equation}
        Rearrange to get common differential equation solution
%        \newline \textcolor{red}{How do I get common solution}
        \begin{equation} \label{eq: a2-8}
            \frac{d^2Q(\phi)}{d\phi^2} = -m^2 Q(\phi)
        \end{equation}
        Common solutions if $m\neq0$
        \begin{equation} \label{eq: a2-9}
            Q(\phi) = A_m e^{im\phi} + B_m e^{-im\phi}
        \end{equation}
        Common solutions if $m=0$
        \begin{equation} \label{eq: a2-10}
            Q(\phi) = A_m + B_m \phi
        \end{equation}
        Back to Original equation \ref{eq: a2-7}, multiply $sin^2(\theta)$ 
        \begin{equation} \label{eq: a2-11}
            \frac{r^2}{R(r)} \frac{d^2 R(r)}{dr^2} + \frac{1}{sin(\theta) P(\theta)} \frac{d}{d\theta}(sin(\theta) \frac{dP(\theta)}{d\theta}) - \frac{m^2}{sin^2\theta)} = 0
        \end{equation}
        First and Last terms are independent so they can be set to constants
        \begin{equation} \label{eq: a2-12}
            \frac{r^2}{R(r)} \frac{d^2 R(r)}{dr^2} - l (l+1) = 0
        \end{equation}
        Where 
        \begin{equation} \label{eq: a2-13}
            -l(l+1) = \frac{1}{sin(\theta) P(\theta)} \frac{d}{d\theta}(sin(\theta) \frac{dP(\theta)}{d\theta}) - \frac{m^2}{sin^2\theta)}
        \end{equation}
        Put equation \ref{eq: a2-12} into more intuitive forms
        \begin{equation} \label{eq: a2-14}
            \frac{d^2 R(r)}{dr^2} = l(l+1) \frac{R(r)}{r^2} 
        \end{equation}
        Put equation \ref{eq: a2-13} into more intuitive forms
        \begin{equation} \label{eq: a2-15}
            \frac{d}{d\theta} (sin(\theta) \frac{d P(\theta)}{d\theta}) + (l(l+1)-\frac{m^2}{sin^2\theta)}) sin(\theta) P(\theta) = 0
        \end{equation}
        Reduce equation \ref{eq: a2-15} to simpler form by substituting $x=cos(\theta)$ and $\frac{d}{d\theta} = -\sqrt{1-x} \frac{d}{dx}$
        \begin{equation} \label{eq: a2-16}
            \frac{d}{dx} ((1-x^2) \frac{d P(X)}{dx}) + (l(l+1)-\frac{m^2}{1-x^2}) P(x) = 0
        \end{equation}
        Equation \ref{eq: a2-13} by substituting $R(r) = r^{\alpha}$
        \begin{equation} \label{eq: a2-17}
            \alpha (\alpha - 1) r^{\alpha - 2} = (l+1) r^{\alpha - 2}
        \end{equation}    
%        \textcolor{red}{Do in between}
        \begin{equation} \label{eq: a2-18}
            \alpha^2 \alpha - 1) - l(l+1) = 0
        \end{equation}        
%        \textcolor{red}{Do in between}
        \begin{gather*}
            \alpha = l + 1
            \\ \alpha = -l
        \end{gather*}
        General solution is 
%        \newline \textcolor{red}{Do in between}
        \begin{equation} \label{eq: a2-19}
            R(r) = A_l r^{l+1} + B_l r^{-1}
        \end{equation} 
        General Solution if the whole azimuthal sweep is included
        \begin{equation} \label{eq: a2-20}
            \Phi(r,\theta, \phi) = \sum_m \sum_l (A_l r^l + B_l r^{-l-1}) (A_m d^{im \phi} + B_m e^{-im \phi} P_l^m(cos(\theta))
        \end{equation}     
        Solving for $P(\theta)$ part we consider a specific case ($m=0$) into equation \ref{eq: a2-16}
        \begin{equation} \label{eq: a2-21}
            \frac{d}{dx} ((1-x^2) \frac{d P(X)}{dx}) + l(l+1) P(x) = 0
        \end{equation}
        Guess, or "Try", power series solution 
        \begin{equation} \label{eq: a2-22}
            P(x) = \sum_{j=0}^{\infty} a_j x^{j+\alpha}
        \end{equation}    
%        \textcolor{red}{Do Inbetween}
        \begin{equation} \label{eq: a2-23}
            \sum_{j=0}^{\infty} (j+\alpha) (j + \alpha - 1) a_j x^{j+\alpha-2} - (j+\alpha) (j + \alpha + 1) a_j x^{j+\alpha} + l(l+1) a_j x^{j+\alpha} = 0
        \end{equation}        
        Consolidate
        \begin{equation} \label{eq: a2-24}
            \sum_{j=0}^{\infty} (j+\alpha) (j + \alpha - 1) a_j x^{j+\alpha-2} - \sum_{j=0}^{\infty} ((j+\alpha) (j + \alpha + 1) + l(l+1)) a_j x^{j+\alpha} = 0
        \end{equation}            
        Remove first two powers of x then combine
%        \newline \textcolor{red}{Do Inbetween}
        \begin{equation} \label{eq: a2-25}
            (\alpha) (\alpha - 1) a_0 x^{\alpha -2) + (1 + \alpha) (\alpha) a_1 x^{\alpha - 2} \sum_{j=0}^{\infty} ((j + 2 + \alpha) (j + \alpha + 1) a_{j + 2} - ((j + \alpha) (j + \alpha + 1) - ((j + \alpha) (j + \alpha + 1) a_j) x^{j+ \alpha}
        \end{equation}            
        This must hold for all values of x so
        \begin{equation} \label{eq: a2-26}
            (\alpha) (\alpha - 1) a_0 = 0
        \end{equation}            
        \begin{equation} \label{eq: a2-27}
            (1 + \alpha) (\alpha) a_1 = 0
        \end{equation}            
        Substitute equation
%        \newline \textcolor{red}{\ref{} into \ref{}}
%        \newline \textcolor{red}{Do Inbetween}
        \begin{equation} \label{eq: a2-28}
            (j + \alpha + 2) (j + \alpha + 1) a_{j + 2} - ((j + \alpha) (j + \alpha + 1 - l(l+1))) = 0
        \end{equation}  
        The third equation is solved to yield the recurrence relation which gives all the remaining terms from the last one. Rerrange equation
%        \newline \textcolor{red}{\ref{}}
        \begin{equation} \label{eq: a2-29}
            a_{j+2} = \frac{(j+\alpha)(j+\alpha+1) - l(l+1)}{(j + \alpha + 2) (j + \alpha + 1} a_j
        \end{equation}      
        $a_1 = 0$; therefore, $a_{odd} = 0$
        \begin{equation} \label{eq: a2-30}
            P_l(x) = \sum_{j=0,even}^{\infty} a_j x^{j+\alpha}
        \end{equation}  
        Convergest at $x=1$
        \begin{equation} \label{eq: a2-31}
            \frac{(j_{max}+\alpha)(j_{max}+\alpha+1) - l(l+1)}{(j_{max} + \alpha + 2) (j_{max} + \alpha + 1} =0
        \end{equation}      
        If $a=0$
        \begin{equation} \label{eq: a2-32}
            j_{max} (j_{max} + 1) - l(l+1) =0
        \end{equation}      
        Therefore 
        \begin{equation} \label{eq: a2-33}
            j_{max}=l
        \end{equation}          
        Substitute equation
%        \newline \textcolor{red}{\ref{} with \ref{}}
        \begin{equation} \label{eq: a2-34}
            P_l(x) = \sum_{j=0,even}^{l} a_j x^{j}
        \end{equation}      
        Therefore 
        \begin{equation} \label{eq: a2-35}
            P_l(x) = \sum_{j=0}^{\frac{l}{2}} a_{2j }x^{2j}
        \end{equation}    
        Rewrite sums over all integers by doubling indices. If l is even
        \begin{equation} \label{eq: a2-36}
            a_{j+2} = j \frac{(j+1) - l(l+1)}{(j + 2) (j + 1)} a_j
        \end{equation} 
        If $\alpha = 1$
        \begin{equation} \label{eq: a2-37}
            (j_{max} + 2)(j_{max} + 1) - l(l+1) =0
        \end{equation} 
        Therefore
        \begin{equation} \label{eq: a2-38}
            j_{max}=l-1
        \end{equation} 
        Substitute equation
%        \newline \textcolor{red}{\ref{} with \ref{}}
        \begin{equation} \label{eq: a2-39}
            P_l(x) = \sum_{j=0,even}^{l-1} a_j x^{j+1}
        \end{equation} 
        Rewrite sums over all integers by doubling indices. If l is odd
        \begin{equation} \label{eq: a2-40}
            a_{j+2} = \frac{(j+1) (j+2) - l(l+1)}{(j + 3) (j + 2)} a_j
        \end{equation}     
        Explicitly state low values of $l$
        \begin{gather*}
            P_0(x) = a_0
            \\ P_1(x) = a_0x
            \\ P_2(x) = a_0 + a_2 x^2 = a_0(1-3x^2) \text{ where $a_2 = -3a_0$}
            \\ P_3(x) = a_0x + a_2x^3 = a_0(x-\frac{5}{3}x^3) \text{ where $a_2 = -\frac{5}{3} a_0$}
        \end{gather*}
        $a_0$ is a scale factor and it is conventional to set $P_l(1)=1$. This conditions allows us to set the scale factor. Additionally, $x$ is a placeholder for the more functional $cos(\theta)$
        \begin{gather*}
            P_0(cos(\theta)) = 1
            \\ P_1(cos(\theta)) = cos(\theta)
            \\ P_2(cos(\theta)) = \frac{1}{2}(-1+3cos^2(\theta))
            \\ P_3(cos(\theta)) = \frac{1}{2}(-3x + \frac{5}{3}cos^3(\theta))
        \end{gather*}
        \newline Legendre polynomials have the following mathematical properties
        \begin{gather*}
            \text{Odd or even symmetry about origin} P_l(-x) = (-1)^l P_l(x) 
            \\ \text{Rodrigues' Formula 1} P_l(x) = \frac{1}{2^l l!} \frac{d^l}{dx^l} (x^2 - 1)^l
            \\ \text{Recurrence relation 1} P_{l+1}(x) = \frac{2 l + 1}{l+1} P_l(x) - \frac{l}{(l+1)} P_{l-1}(x)
            \\ \text{Rodrigues' Formula 1} P_l(x) = \frac{1}{2 l + 1} \frac{d}{dx} (P_{l+1}(x) - P_{l-1}(x))
            \\ \text{Orthogonality condition} \int_{-l}^l P_{l'}(x) P_l(x)dx = \frac{2}{2l + 1} \delta_{l'l}
            \\ \text{Orthogonality condition (Polar angle)} \int_0^\pi P_{l'}(cos(\theta)) P_l(cos(\theta)) sin(\theta) d\theta = \frac{2}{2l+1}\delta_{l'l}
        \end{gather*}
        The Legendre Polynomials form a complete set of orthogonal functions on the interval $(-1,1)$, so any function $f(x)$ and be expanded in terms of Legendre Polynomials
        \begin{equation} \label{eq: a2-41}
            f(x) = \sum_{l=0}^\infty A_l P_l(x)
        \end{equation}
        Where 
        \begin{equation} \label{eq: a2-42}
            A_l = \frac{2l+1}{2} \int_{-1}^1 f(x) P_l(x) dx
        \end{equation}
        The general solution to the Laplace Equation in spherical coordinates for the case $m=0$ (where the origin and non-origin are considered)
        \begin{equation} \label{eq: a2-43}
            \Phi(r,\theta, \phi) = \sum_{l=0}^\infty (A_l r^l + B_l r^{-l-1}) P_l(cos(\theta))
        \end{equation}
    
\newpage
\printbibliography
\end{document}
